\section{Conclusion}
\label{sec:conclusion}

This paper presented an integrated, two-stage framework for lightweight ECG denoising and automated cardiac arrhythmia classification. In Stage~1, we introduced \HSL, a 75{,}133-parameter one-dimensional wavelet residual network combining discrete Haar multiresolution analysis, unified standard residual blocks with SE1D channel attention (\texttt{ResSEBlock1D}), detail subband residual filtering, and a physics-preserving mixed loss objective. On the MIT-BIH Arrhythmia Database corrupted with NSTDB noise across 7 noise categories and 10 SNR tiers, \HSL{} established superior reconstruction fidelity, achieving an average $\Delta\mathrm{SNR}$ of \best{10.01\,dB} and a $\mathrm{PRD}$ of \best{33.88\%}, while executing in only \best{35.20\,ms} per segment on an edge CPU.

In Stage~2, we proposed \RRN, coupling the frozen denoiser to a 1D ResNet morphology encoder and an independent 6-dimensional RR contextual pathway. Evaluated on the gold-standard de Chazal inter-patient DS1/DS2 benchmark (49{,}298 independent test beats), 5-seed soft-voting ensembles achieved \best{96.43\%} accuracy and \best{0.8407} Macro-F1 with a classical MLP head (boosting S-class precision to \best{72.97\%}), and \best{96.04\%} accuracy and \best{0.8176} Macro-F1 with a 6-qubit hybrid variational quantum head. Ablation analyses proved that channel attention (SE1D) preserves high-frequency QRS fiducials substantially better than spatial attention (CBAM1D). By bridging the divide between waveform enhancement and diagnostic accuracy, this research delivers an ultra-compact, clinically validated foundation for next-generation wearable cardiovascular telemetry.
